\chapter*{后记: 接下来往哪里走?}\label{chap:coda}
\markboth{后记}{接下来往哪里走?}
\setcounter{equation}{0}
\renewcommand{\theequation}{C.\arabic{equation}}
\renewcommand{\theHequation}{coda.\arabic{equation}}

本书已经讲了大量的数学, 但说我们不过刚刚触及代数几何的表面, 并不夸张. 这本书是一道门. 如果读者一路认真读了下来, 就已走进代数几何这座房子, 准备好去探索其中形形色色的房间. 不过, 先选择哪个房间探索, 在某种意义上取决于个人兴趣和机缘. 在这篇后记中, 我们介绍几个彼此补充的方向, 供读者读完本书后继续学习. 前两个方向聚焦于两类性质特别好, 研究得特别透彻的簇, 即代数曲线和环面簇; 第三个方向则是彻底推广簇的概念, 进入概形论的广阔天地. 希望这篇后记能让读者初步体会前方还有什么, 并为进一步掌握代数几何这门浩瀚而优美的学科指明一些方向.

\section*{代数曲线}

学习射影代数几何的许多工具, 或许最好先深入研究一维射影簇, 也就是代数曲线. 更具体些, 可以先研究射影平面曲线, 即由 $\mathbb{P}^{2}$ 中单个多项式的消没定义的射影簇. 这个范围看似很窄, 却已是一个丰富而迷人的世界, 其中包括对交理论工具的极好介绍.

关于射影平面曲线的交理论断言, 最典型的例子是本书曾简要提到的 Bézout 定理. 粗略地说, Bézout 定理断言: 若 $C_{1}=\mathcal{V}(f_{1})$ 和 $C_{2}=\mathcal{V}(f_{2})$ 是没有公共不可约分支的射影平面曲线, 且 $\deg(f_{i})=d_{i}$, 则 $C_{1}$ 与 $C_{2}$ 交于 $d_{1}\cdot d_{2}$ 个点. 然而, 涉及相切时, 必须对这句话加以限定. 例如, 曲线 $C_{1}=\mathcal{V}(y)$ 和 $C_{2}=\mathcal{V}(yz-x^{2})$ 分别由一次和二次多项式定义, 所以 Bézout 定理似乎应当告诉我们它们有两个交点, 而事实上它们的交只有一个点 $[0:0:1]$. 解决办法是注意到 $C_{1}$ 与 $C_{2}$ 在 $[0:0:1]$ 处相切, 并把这个点看成\textquotedbl{}重数\textquotedbl{}为二的交点. 一般而言, 只有按适当的重数来计数交点, Bézout 定理才成立. 精确定义交重数, 并证明它满足 Bézout 定理这样的性质, 正是交理论的起点. 对于希望在代数曲线情形下学习交重数和 Bezout 定理的学生, Fulton 的书~[\ref{bib:fulton}] 是一本经典参考书.

除了可以借此入门交理论, 曲线还有另一处值得关注: 当基域是 $\mathbb{C}$ 时, 可以借助复分析研究代数几何. 关键在于, $\mathbb{C}$ 上的光滑一维射影簇可以视为一维紧致复流形, 也称为 Riemann 曲面. 任意维的光滑射影簇也有类似性质, 但在一维时, 奇妙的是逆命题也成立: 每个紧致 Riemann 曲面作为复流形都同构于一个光滑一维射影簇. 因此, 既可以用代数几何的方法研究曲线 (把它们看成射影簇), 也可以用复分析的方法研究曲线 (把它们看成 Riemann 曲面). 代数观点和分析观点之间有着优美的相互作用. 如果读者觉得以分析的眼光看待曲线颇有启发, 可以学习 Cavalieri 和 Miles 的书~[\ref{bib:cavalieri-miles}] 或 Miranda 的书~[\ref{bib:miranda}].

除了几何以外, 代数曲线还与数论有丰富而重要的联系. 为略见一斑, 回顾鼎鼎有名的\textquotedbl{}Fermat 最后定理\textquotedbl{}: 对任意 $n\ge3$, 不存在使 $x^{n}+y^{n}=z^{n}$ 成立的正整数 $x,y,z$. 用代数几何的话说, 这个断言就是: 当 $n\ge3$ 时, 射影平面曲线 $\mathcal{V}(x^{n}+y^{n}-z^{n})\subseteq\mathbb{P}_{\mathbb{C}}^{2}$ 上不存在坐标均为非零有理数的点. 这只是数论与代数曲线上有理点的研究密切相关的初步例证. 若想从数论观点入门曲线上有理点的理论, 读者可以学习 Silverman 和 Tate 的书~[\ref{bib:silverman-tate}].

\section*{环面簇}

另一个可供继续学习的方向是环面几何, 它是进入高于一维的代数几何的一条具体且便于动手实践的途径. 简单说, 环面簇是 $\mathbb{C}$ 上的代数簇 $X$, 它含有一个与\textquotedbl{}代数环面\textquotedbl{} $T\cong(\mathbb{C}^{*})^{n}$ 同构的稠密开集, 且 $T$ 通过逐分量乘法在自身上的群作用可以延拓为 $T$ 在整个 $X$ 上的作用. 例如, 射影空间 $\mathbb{P}^{n}$ 就是环面簇, 其稠密开集由齐次坐标均非零的点构成.

环面几何的核心思想是, 只要 $X$ 上存在这个稠密的 $T$ 作用, $X$ 就具有丰富的组合结构, 从而可以完全用离散几何的语言重新解释 $X$ 的几何. 例如, 射影簇 $X$ 的一项重要代数几何信息是它的次数, 它通过计数 $X$ 与一互补维数的线性空间的交点, 粗略地衡量 $X$ 有多\textquotedbl{}弯曲\textquotedbl{}. 当 $X$ 是环面簇时, 它的次数等于相应多胞体 $P_{X}$ 的体积. 这样, 环面几何把代数几何中的一个量 (次数) 转化成了离散几何中的一个量 (体积).

凭借这种向离散情形的转化, 环面几何可以让学生在具体且便于演算例子的环境中, 初步体会代数几何中许多常用工具, 从除子理论和奇点消解, 到上同调和交理论. 不过, 须事先提醒, 在一般代数簇的广阔天地中, 环面簇只占很小一部分; 例如, 唯一的一维光滑射影环面簇是 $\mathbb{P}^{1}$. 有兴趣进一步学习环面几何的学生, 可以查阅 Cox, Little 和 Schenck 的详尽教材~[\ref{bib:cox-little-schenck}].

\section*{概形论}

这里最后提到的一个深造主题范围极广, 本身就是通往现代代数几何其他各部分的一道门, 那就是概形.

为引出概形的想法, 回顾仿射簇与有限生成既约 $K$-代数之间有一一对应 $X\mapsto K[X]$. 理解其逆对应的一种办法是回顾, $X$ 的点与 $K[X]$ 的极大理想一一对应, 所以至少作为集合, 可以把 $X$ 恢复为极大谱
\[
  \operatorname{maxSpec}(K[X])=\{\text{极大理想 }M\subseteq K[X]\}.
\]
事实上, 借此不仅可以把 $X$ 恢复为集合, 还可以恢复为拓扑空间: 仿射簇 $X$ 的 Zariski 闭子集正是形如 $\mathcal{V}(I)\subseteq X$ 的集合, 其中 $I\subseteq K[X]$ 是理想. 按照它与 $\operatorname{maxSpec}(K[X])$ 的一一对应逐项解释, 就在 $\operatorname{maxSpec}(K[X])$ 上得到一个 Zariski 拓扑, 其闭集恰为任意理想 $I\subseteq K[X]$ 所对应的集合 $Z_{I}=\{\text{极大理想 }M\supseteq I\}$.

这样, 我们描述了一个单射
\begin{equation}\label{eq:coda-maximal-spectrum}
  \begin{aligned}
    \{\text{有限生成既约 }K\text{-代数}\}&\to\{\text{拓扑空间}\}\\
    R&\mapsto\operatorname{maxSpec}(R)
  \end{aligned}
\end{equation}
其像恰由带 Zariski 拓扑的仿射簇组成. 然而, 左边是相当特殊的一类环. 因此, 自然要问, 能否将 \eqref{eq:coda-maximal-spectrum} 扩张为任意交换环 $R$ 与仿射簇的某种推广之间的一一对应. 前文已略微提示为何这样的推广会有用: 仿射曲线 $\mathcal{V}(y)\subseteq\mathbb{A}^{2}$ 与 $\mathcal{V}(y-x^{2})\subseteq\mathbb{A}^{2}$ (本篇后记前面所考察射影平面曲线的仿射限制) 相交于
\[
  X=\mathcal{V}(y)\cap\mathcal{V}(y-x^{2})=\mathcal{V}(y,y-x^{2}).
\]
若忘记 $\langle y,y-x^{2}\rangle$ 不是根理想, 就会误以为 $K[X]$ 是环
\[
  \frac{K[x,y]}{\langle y,y-x^{2}\rangle}\cong\frac{K[x]}{\langle x^{2}\rangle},
\]
它不是既约环. 这个环不是既约的, 特别是它由一个平方为零的元素生成, 这恰好反映了 $\mathcal{V}(y)$ 与 $\mathcal{V}(y-x^{2})$ 以\textquotedbl{}重数二\textquotedbl{}相交. 因此, 虽然 $K[x,y]/\langle y,y-x^{2}\rangle$ 不是簇 $X$ 的坐标环, 却可以尝试将它视为对 $X$ 的某种加强的坐标环, 这种加强记录了 $X$ 原本是一个切点这一事实.

那么, 或许只需把 \eqref{eq:coda-maximal-spectrum} 中有限生成既约 $K$-代数换成任意交换环. 遗憾的是, 这不可能是正确答案: 对任意交换环, 对应 $R\mapsto\operatorname{maxSpec}(R)$ 不是单射. 作为说明这一问题的例子, 请读者验证 $K[x]/\langle x\rangle$ 和 $K[x]/\langle x^{2}\rangle$ 各只有一个极大理想, 但两者并非同构的环.

解决办法是考虑比极大谱略大且性质更好的集合, 即素谱
\[
  \operatorname{Spec}(R)=\{\text{素理想 }P\subseteq R\}.
\]
完全按照给极大谱赋予 Zariski 拓扑的方式, 可以给这个集合赋予 Zariski 拓扑. 此外, $\operatorname{Spec}(R)$ 还从环 $R$ 继承了额外的结构, 称为\textquotedbl{}结构层\textquotedbl{} (下面再细述). 一个拓扑空间连同其上的一个环层称为\emph{赋环空间}. 由此得到映射
\begin{equation}\label{eq:coda-prime-spectrum}
  \begin{aligned}
    \{\text{交换环}\}&\to\{\text{赋环空间}\}\\
    R&\mapsto\operatorname{Spec}(R),
  \end{aligned}
\end{equation}
它终于是一个单射. 按定义, 仿射概形就是此映射的像中的赋环空间, 即形如 $\operatorname{Spec}(R)$ 的集合连同其 Zariski 拓扑和结构层, 其中 $R$ 为某个交换环. 概形论的出发点是, 将 \eqref{eq:coda-prime-spectrum} 的陪域限制为仿射概形后, 它就成为一个范畴等价, 极大地推广了贯穿本书的代数与几何的等价.

到目前为止, 对概形的讨论完全在仿射情形下进行, 于是自然会问: 是否有一个概念, 它与 (拟) 射影簇的关系, 如同仿射概形与仿射簇的关系? 这需要细致得多的处理, 因为仿射簇 $X$ 虽然完全由其正则函数环 $K[X]$ 决定, 拟射影簇却不满足这一结论.

实际上, 正如 Section~\ref{sec:continuity-of-regular-maps} 末尾所提示的, 决定拟射影簇 $X$ 的是所有开集 $U\subseteq X$ 上的环 $K[U]$ 的数据. 当两个开集 $U,V$ 相交时, 它们通过限制映射 $K[U]\to K[U\cap V]$ 联系起来. 这些数据本质上就是所谓 $X$ 上的环层. 更精确一点说, 任意拓扑空间 $X$ 上的一个\emph{环层}为每个开集 $U\subseteq X$ 指定一个环 $R(U)$, 并在 $U'\subseteq U$ 时指定一个环同态 $R(U)\to R(U')$, 这些指定须满足某些相容性条件.

特别地, 若 $X$ 是拟射影簇 (视为带有 Zariski 拓扑的拓扑空间), 则 $X$ 上有一个称为\emph{结构层}的环层 $\mathcal{O}_{X}$, 对任意开集 $U\subseteq X$ 定义为
\[
  \mathcal{O}_{X}(U)=K[U].
\]
事实上, 拟射影簇 $X$ 由赋环空间 $(X,\mathcal{O}_{X})$ 决定, 所以在这个意义下, 拟射影簇可视为赋环空间的特例. 不过, 并非所有赋环空间都是拟射影簇. 赋环空间 $(X,\mathcal{O}_{X})$ 的一个重要特征是\textquotedbl{}局部仿射\textquotedbl{}, 意思是它有一个由开集 $U_{\alpha}\subseteq X$ 组成的开覆盖, 对每个 $U_{\alpha}$, 限制 $(U_{\alpha},\mathcal{O}_{X}|_{U_{\alpha}})$ 都同构于一个仿射簇连同其结构层. 满足此条件的赋环空间 (不一定是拟射影簇) 称为\emph{抽象代数簇}. 用层的语言重述簇的概念以后, 就可以继续把概形推广到仿射情形之外: 首先, 如上所述, 描述一种在任意交换环 $R$ 的仿射概形 $\operatorname{Spec}(R)$ 上自然构造结构层的方法; 然后, 定义概形为一个赋环空间, 它有一个由作为赋环空间同构于仿射概形的开集构成的开覆盖.

概形的语言初看可能令人畏惧, 但任何希望进一步学习代数几何的学生都必须努力掌握它. 建议读者查阅 Eisenbud 和 Harris 的书~[\ref{bib:eisenbud-harris}], Hartshorne 的经典著作~[\ref{bib:hartshorne}] 或 Vakil 的书~[\ref{bib:vakil}], 开始把这些概念收入自己的数学工具箱.

\par\medskip
{\footnotesize
\noindent © 作者 2025\par
\noindent E.\ Clader, D.\ Ross, \emph{Beginning in Algebraic Geometry}, Undergraduate Texts in Mathematics.\par
\noindent
\href{https://doi.org/10.1007/978-3-031-88819-9}{doi:10.1007/978-3-031-88819-9}.\par
}
